%%%% APÊNDICE A
%%
%% Texto ou documento elaborado pelo autor, a fim de complementar sua argumentação, sem prejuízo da unidade nuclear do trabalho.

%% Título e rótulo de apêndice (rótulos não devem conter caracteres especiais, acentuados ou cedilha)

\definecolor{lightgray}{gray}{0.95}

\lstset{
    backgroundcolor=\color{lightgray},
    basicstyle=\ttfamily\small,
    breaklines=true,
    numbers=left,
    numberstyle=\tiny,
    frame=single,
    tabsize=2
}

\chapter{Algoritmo de Coleta e Configuração do Servidor}\label{cap:apendicea}
\begin{lstlisting}[language=bash, caption={Configuração do Nginx utilizada nos experimentos}, label={lst:nginx-config}]
user  nginx;

-- MANTENDO O EXEMPLO

    include /etc/nginx/conf.d/*.conf;
}
\end{lstlisting}

\newpage

\begin{lstlisting}[language=bash, caption={Configuração do servidor Nginx com HTTP/2 e HTTP/3 (QUIC)}, label={lst:nginx-http3-config}]
server {
    listen 443 ssl;
MANTENDO O EXEMPLO
\end{lstlisting}

\newpage

%% Título e rótulo de seção (rótulos não devem conter caracteres especiais, acentuados ou cedilha)
%\section{Título da Seção Secundária do Apêndice B}\label{sec:secaoapendicea}

%Exemplo de seção secundária em apêndice (\autoref{sec:secaoapendicea} do \autoref{cap:apendicea}).

%% Título e rótulo de seção (rótulos não devem conter caracteres especiais, acentuados ou cedilha)
%\subsection{Título da Seção Terciária do Apêndice B}\label{subsec:subsecaoapendicea}

%Exemplo de seção terciária em apêndice (\autoref{subsec:subsecaoapendicea} do \autoref{cap:apendicea}).

%% Título e rótulo de seção (rótulos não devem conter caracteres especiais, acentuados ou cedilha)
%\subsubsection{Título da seção quaternária do Apêndice B}\label{subsubsec:subsubsecaoapendicea}

%Exemplo de seção quaternária em apêndice (\autoref{subsubsec:subsubsecaoapendicea} do \autoref{cap:apendicea}).

%% Título e rótulo de seção (rótulos não devem conter caracteres especiais, acentuados ou cedilha)
%\paragraph{Título da seção quinária do Apêndice B}\label{para:paragraphapendicea}

%Exemplo de seção quinária em apêndice (\autoref{para:paragraphapendicea} do \autoref{cap:apendicea}).
